\section{Paired Information Revealing: Formal Setup}
\label{sec:setup}

We formalize paired information revealing. Let
$q:\mathcal{X}\rightarrow\mathbb{R}$ denote the model score.
An input $\mathbf{x}^{+}$ and a factor-level counterfactual
$\mathbf{x}^{-}$ define the final contrast
\[
d=q(\mathbf{x}^{+})-q(\mathbf{x}^{-}).
\]
The pair need not be pre-specified: in feature attribution,
$\mathbf{x}^{-}$ can be induced from $\mathbf{x}^{+}$ by a
removal or replacement operator.

A paired reveal starts from a common uninformative state $\mathbf{x}_0$
and produces matched trajectories
$(\mathbf{x}^{+}(t),\mathbf{x}^{-}(t))$ for $t\in \mathcal{T}=[0,1]$, with
\[
\mathbf{x}^{+}(0)=\mathbf{x}^{-}(0)=\mathbf{x}_0,\qquad
\mathbf{x}^{+}(1)=\mathbf{x}^{+},\qquad
\mathbf{x}^{-}(1)=\mathbf{x}^{-}.
\]

The signed response is
$r(t)=q(\mathbf{x}^{+}(t))-q(\mathbf{x}^{-}(t))$, with $d=r(1)$.
Ordinary magnitude $|r(t)|$ preserves size but discards its
relation to the final contrast; Section~\ref{sec:method} uses it
to assign semantic roles.